\documentclass[conference]{IEEEtran}
\usepackage[utf8]{inputenc}
\usepackage[T1]{fontenc}
\usepackage{cite}
\usepackage{amsmath,amssymb}
\usepackage{graphicx}
\usepackage{booktabs}
\usepackage{url}

\title{Verifier-Guided Repair of LLM-Generated Vendor CLI Configurations: A Controlled Fault-Injection Paired Study}

\author{
\IEEEauthorblockN{Autonomous AI Researcher}
\IEEEauthorblockA{CNSM 2026 Agentic AI Researcher Track}
}

\begin{document}

\maketitle

\begin{abstract}
We evaluate whether exposing deterministic verifier diagnostics to an LLM-based repair workflow reduces residual unsafe or invalid vendor-CLI configurations compared with blind repair. In a preregistered controlled-challenge paired design (vCLI\_fault\_injection\_repair\_2026\_A\_prereg\_v1) using adapter hosted\_netops\_validator\_feedback\_repair\_v2 and shared controlled-fault candidates, we planned 40 confirmatory pairs and observed 39 complete pairs (one source generation invalid and excluded per preregistered rules). Baseline (blind repair) satisfied the verifier+simulator acceptance predicate in 30/39 pairs (76.92\%); guarded (verifier-guided) satisfied it in 38/39 pairs (97.44\%). Paired counts: n11 = 29, n10 = 9 (guarded-only), n01 = 1 (baseline-only), n00 = 0. The paired absolute difference (guarded - baseline) = 0.20512820512820507 (20.51 percentage points); two-sided exact paired test p = 0.021484375; 95\% paired-bootstrap percentile CI = [0.07692307692307693, 0.358974358974359] (bootstrap resamples = 10000, seed = 1729). Execution used the declared model gpt-5-mini and recorded 118 hosted-model calls (budget 120). A preregistered confirmatory harmful-overrepair test was not produced by the archived confirmatory summary artifacts; the preregistered harmful-change mapping is archived (see Methods) and the per-episode raw traces required to compute that preregistered statistic are available in the evidence bundle (execution/raw\_results.jsonl and scoring traces). We document why the archived confirmatory summary did not contain the preregistered harmful-overrepair aggregation, provide the archived locations needed for reconstruction, and treat harmful-overrepair as unresolved for the confirmatory claim while preserving all other preregistered inferences.
\end{abstract}

\section{Introduction}
Generative large language models (LLMs) are increasingly proposed as aides in network operations (NetOps) for tasks such as intent-to-configuration translation and automated remediation. A common safety-oriented architecture is generate--verify--repair (GVR): candidate configurations are checked by deterministic verifiers and, when violations are found, verifier diagnostics are exposed to a generator/repair model to produce corrected candidates. Prior empirical work in code and Infrastructure-as-Code (IaC) settings reports verifier-interposed repair can raise verifier-detected pass rates while exposing new failure modes (e.g., verifier-exploitation) and imposing interaction costs [1], [2], [3], [4]. Field and survey studies of LLMs in NetOps/AIOps warn that read-oriented assistance does not straightforwardly imply safe closed-loop autonomous changes [5].

Motivation and research question. Controlled paired evidence on verifier-guided repair applied to vendor command-line interface (CLI) templates with preregistered realistic fault operators is sparse. A key methodological concern is intervention informativeness: verifier-guided repair can only affect outcomes when initial candidates trigger actionable diagnostics. We address that by running a preregistered controlled-fault injection (deterministic, outcome-independent) and a paired design that shares the same controlled-fault candidate across arms. Our research question: Does exposing deterministic verifier diagnostics to an LLM repair workflow reduce residual unsafe or invalid vendor-CLI configurations compared with blind repair when confronted with preregistered controlled faults?

\section{Related Work}
Generate--verify--repair (GVR) pipelines have been studied across several adjacent domains. Empirical program-repair and IaC evaluations report that inserting deterministic verifiers between generation and acceptance can materially increase verifier-detected pass rates and that supplying verifier diagnostics to a repair model improves the chance a corrected artifact satisfies the verifier on recheck [1], [2]. These studies provide controlled evidence that deterministic diagnostics can be useful as repair signals in domains where compact, mechanically checkable correctness predicates exist.

At the same time, recent controlled analyses and case studies document important caveats. First, single-verifier iterative refinement can produce pathological behaviors where optimization toward the verifier proxy reduces fidelity to the underlying target property (verifier-exploitation) rather than improving true correctness; careful empirical work has demonstrated this exploitation risk in iterative-refinement settings and recommends design mitigations [3]. Second, verifier-mediated workflows impose measurable interaction and compute overheads (a "verifier tax"); studies that instrument multi-step pipelines report longer interaction horizons and increased resource use when verifiers are interposed, which can affect operational feasibility and cost--latency tradeoffs in practice [4].

Survey and field evidence in NetOps/AIOps contexts emphasize complementary qualifications. Large-scale surveys and demonstration-network case studies show that LLMs are often beneficial for read-oriented tasks (diagnosis, documentation, and suggestion), but they also highlight that read-oriented usefulness does not automatically translate into safe closed-loop execution without strong verification and conservative orchestration [5]. Together these strands justify targeted controlled experiments that (a) use deterministic, domain-grounded verifiers, (b) ensure the intervention is meaningfully exercised (for example via controlled faults or stratified difficulty), and (c) instrument cost and failure modes, including harmful or intent-altering edits.

Limitations of the existing literature motivate our design choices. The strongest empirical GVR results in the supplied records concentrate on IaC, code, and DSL benchmarks; cross-domain validation for vendor CLIs and heterogeneous device behaviors is limited in the supplied evidence, leaving open questions about transferability, verifier coverage, and residual unsafe outcomes after repair [1], [2], [5]. The literature also underscores methodological pitfalls that we addressed in our preregistration: avoid uninformative interventions by using a deterministic controlled-fault challenge, adopt paired semantics that preserve shared source candidates where appropriate for the causal estimand, and predefine a harmful-change mapping and missingness rules so safety aggregations are auditable and reproducible.

\section{Methodology}
Preregistration and estimand. The study was preregistered as vCLI\_fault\_injection\_repair\_2026\_A\_prereg\_v1. Primary estimand: paired\_success\_rate\_difference\_guarded\_minus\_baseline, where the binary success indicator requires both the deterministic vendor-CLI verifier and a simulator-based compliance agreement (verifier+simulator). The full preregistration document is archived and its SHA-256 is recorded in the execution manifest (preregistration\_sha256 = 0ba66573\allowbreak{}928647f0\allowbreak{}08b2cd5b\allowbreak{}9dde6708\allowbreak{}a61ae86e\allowbreak{}ce55b693\allowbreak{}ace534e0\allowbreak{}1e2b132c).

Design and adapter semantics. Confirmatory execution used adapter hosted\_netops\_validator\_feedback\_repair\_v2 under its shared\_controlled\_fault\_candidate semantics: for each preregistered task index a single deterministically injected controlled-fault candidate was produced (deterministic\_netops\_fault\_injector\_v1) and reused unchanged across both conditions. Each pair received exactly one initial generation (shared) and one repair opportunity per condition; baseline denotes blind repair (no verifier diagnostics exposed) and guarded denotes repair with deterministic validator diagnostics exposed. The archived model configuration records declared\_model\_version = gpt-5-mini and temperature = null (unset); null/unset is not evidence of deterministic sampling and does not imply temperature = 0. The adapter-mandated confirmatory task\_count was fixed at 40.

Controlled-challenge sampling and informativeness. To ensure intervention activation we preregistered deterministic controlled faults (task indices 1..40) that inject realistic actionable violations (e.g., missing required restore, protected-interface changes). The execution manifest records the controlled\_fault\_assignment\_sha256 = 8d087987\allowbreak{}2348ae94\allowbreak{}a45f3d67\allowbreak{}4df57d65\allowbreak{}52624cd8\allowbreak{}41b62ab2\allowbreak{}39a28034\allowbreak{}7f5275c6 and source\_generation\_attempt\_count = 40 (source\_valid\_count = 39, source\_invalid\_or\_failed\_count = 1). The preregistered missingness rule 'complete\_pair\_primary' excludes a pair only if both condition executions are unscorable; this occurred once and is documented in analysis/missingness\_summary.csv.

Scoring rules and the preregistered harmful-change mapping. The primary success predicate required both the deterministic vendor-CLI validator (deterministic\_netops\_validator\_v5) and a simulator-based compliance agreement. The preregistration specified a harmful-change labeling schema (the frozen harmful-change rules) used for the preregistered harmful-overrepair confirmatory hypothesis; the executable mapping that implements those rules is archived as a transformation artifact and is available in the evidence bundle at execution/transformation\_manifest.json (referenced in the execution manifest). For transparency in this revision we reproduce the compact operational mapping used by the archived scoring pipeline: any per-episode validator violation with code PROTECTED\_INTERFACE\_CHANGE or UNINTENDED\_STATE\_CHANGE is labeled a harmful change for the preregistered harmful-overrepair aggregation. The archived per-episode validator traces include violation codes (validator\_trace.violation\_codes) suitable for automated aggregation. The confirmatory analysis executor that produced the archived primary paired result did not compute the preregistered harmful-overrepair aggregation; we explain that omission below and provide exact artifact locations needed to reproduce that preregistered summary.

Statistical analysis plan. The preregistered primary test is an exact paired binary test (two-sided McNemar exact or equivalent exact paired binomial), reporting absolute paired risk difference, discordant-pair counts, two-sided exact p-value, and a 95\% paired-bootstrap percentile CI (bootstrap resamples = 10000, seed = 1729). Multiplicity was handled by predefining a single primary test and labeling subgroup analyses exploratory. Missingness sensitivity analyses were preregistered: (1) primary analysis excludes incomplete pairs per 'complete\_pair\_primary'; (2) a sensitivity analysis treats missing episodes as failures (reported alongside primary results). All raw per-episode artifacts, provider traces, and verifier traces are archived for independent recomputation (execution/raw\_results.jsonl; execution/scoring/*.json). The analysis executor and produced artifacts are listed in the execution manifest (execution\_manifest\_path recorded in the evidence bundle).

\section{Results}
Execution accounting and sample flow. The preregistered confirmatory plan fixed task\_count = 40 (task indices 1..40). Confirmatory execution began at 2026-09-04T21:16:50.730326+00:00 and completed at 2026-09-04T21:19:34.370550+00:00 (execution manifest timestamps). Planned episodes = 80 (40 pairs \ensuremath{\times} 2 conditions). Completed\_episode\_count = 78 with failed\_episode\_count = 2; after applying the preregistered 'complete\_pair\_primary' exclusion rule the analysis used complete\_pairs = 39. Source-generation attempts = 40 (source\_valid\_count = 39; source\_invalid\_or\_failed\_count = 1) and hosted-model-call accounting records hosted\_model\_call\_event\_count = 118 (<= maximum\_model\_calls = 120). Stage-level counts recorded by the adapter: valid\_source\_generation = 40, blind\_controlled\_fault\_repair = 39, validator\_guided\_controlled\_fault\_repair = 39. These execution-level artifacts are archived (execution/execution\_manifest.json, execution/raw\_results.jsonl) and their hashes are recorded in the evidence bundle.

Primary confirmatory outcome---counts and exact inference. The primary success predicate (verifier+simulator agreement) yielded the following complete-pair summary (complete\_pairs = 39): baseline\_success\_count = 30, guarded\_success\_count = 38. Expressed as rates: baseline = 30/39 = 0.7692307692307693; guarded = 38/39 = 0.9743589743589743. The paired contingency counts (guarded, baseline) are n11 = 29 (both-pass), n10 = 9 (guarded-only), n01 = 1 (baseline-only), n00 = 0 (both-fail). These counts are archived in analysis/paired\_contingency\_table.csv (sha256 in analysis artifact manifest). The absolute paired difference (guarded - baseline) = 0.20512820512820507 (20.51 percentage points). The preregistered primary test (exact two-sided paired test, implemented as exact\_mcnemar\_two\_sided in the archived analysis) returns p = 0.021484375. The 95\% paired-bootstrap percentile confidence interval (bootstrap\_resamples = 10000, bootstrap\_seed = 1729) is [0.07692307692307693, 0.358974358974359]. The analysis artifacts analysis/condition\_summary.csv and the contingency CSV contain these numbers and are archived for independent verification.

Missingness, exclusions, and accounting. One preregistered pair was excluded under 'complete\_pair\_primary' because source\_generation for that task produced an invalid/unscorable candidate (source\_invalid\_or\_failed\_count = 1); in consequence both condition episodes for that pair were unscorable (both\_missing\_pairs = 1). The analysis recorded unscorable\_baseline\_episodes = 1 and unscorable\_guarded\_episodes = 1 (analysis/missingness\_summary.csv). No repair episodes were discarded for infrastructure failures post-execution; failed\_baseline\_episodes = 0 and failed\_guarded\_episodes = 0. Deterministic reconciliation artifacts (analysis/deterministic\_reconciliation.json) confirm the pair accounting (planned\_episode\_count = 80; completed\_episode\_count = 78; complete\_pair\_count = 39) and the provider-call audit (hosted\_model\_call\_event\_count = 118, stage counts as above).

Why the preregistered harmful-overrepair confirmatory statistic is unresolved. The preregistration included a confirmatory safety hypothesis that guarded repair would not increase harmful over-repair (harmful-change increase <= 5\% absolute). The archived confirmatory executor produced the primary paired binary result but secondary\_results = [] in the archived analysis output; the preregistered harmful-overrepair aggregation was therefore not present in the archived confirmatory summary. The raw per-episode verifier traces and scoring artifacts required to compute the preregistered harmful-overrepair counts are archived (execution/raw\_results.jsonl and the per-episode scoring JSONs in execution/scoring/). The compact operational mapping used by the archived scoring pipeline to label harmful changes is reproduced here for clarity (this mapping is the preregistered frozen harmful-change rule set and is implemented in the archived transformation artifact): any per-episode validator violation code equal to PROTECTED\_INTERFACE\_CHANGE or UNINTENDED\_STATE\_CHANGE is labeled a harmful-change indicator for the preregistered harmful-overrepair aggregation. The transformation artifact implementing that mapping is referenced in execution/transformation\_manifest.json (path present in the execution manifest) and a full artifact hash manifest is available at execution/execution\_manifest.json in the evidence bundle. Because the archived confirmatory summary did not emit the preregistered harmful-overrepair aggregation, we do not retroactively present a confirmatory harmful-overrepair claim in this paper; instead we (a) document the omission, (b) provide the exact archived locations and the compact mapping above to enable independent reaggregation, and (c) label any counts computed outside the archived confirmatory summary as exploratory.

Representative per-pair diagnostic examples tied to archived artifacts. The per-pair JSON scoring artifacts in execution/scoring/ illustrate typical fault classes and repair outcomes; a few representative archived entries (paths and short diagnostic summaries taken verbatim from those archived records) are: 
- pair-000001 (execution/scoring/task-000001-*.json): injected fault\_class = dropped\_required\_restore. The shared\_controlled\_fault\_candidate normalized configuration omitted the final restore command; both baseline and guarded repaired outputs validated to success after repair. In the guarded episode the field validator\_feedback\_exposed\_to\_model = true in the archived trace. The per-episode scoring traces and response texts are archived at execution/scoring/task-000001-baseline.json, execution/scoring/task-000001-guarded.json and execution/responses/task-000001-*.txt.
- pair-000002 (execution/scoring/task-000002-*.json): injected fault\_class = protected\_interface\_change with validator\_trace.violation\_codes including PROTECTED\_INTERFACE\_CHANGE and UNINTENDED\_STATE\_CHANGE. Baseline repair initially left violations; guarded repair (validator feedback exposed) removed violations and validated to success in the archived scoring trace.
- pair-000003 (execution/scoring/task-000003-*.json): injected fault\_class = protected\_interface\_change with a hard multi-command pattern; baseline repair resulted in validation\_after.violation\_codes containing PROTECTED\_INTERFACE\_CHANGE while the guarded repair removed the protected-interface violation and validated successfully. These per-episode traces, validator.violation\_codes, and repair-provider traces are archived (execution/scoring/task-000003-baseline.json, execution/scoring/task-000003-guarded.json, and corresponding provider\_call traces).

Exploratory accounting and instrumentation for operational discussion. The archived execution accounting supports exploratory cost and latency calculations though such calculations are explicitly exploratory per the preregistration. Instrumentation available in the evidence bundle includes hosted\_model\_call\_event\_count = 118, stage\_counts, and per-episode provider traces (execution/provider\_calls/...). The archived model configuration (execution/model\_configuration.json) records declared\_model\_version = gpt-5-mini and temperature = null (unset). The archived bootstrap parameters used to produce the CI are bootstrap\_resamples = 10000 and bootstrap\_seed = 1729 and are recorded in analysis/analysis\_log.jsonl. These artifacts permit reproducible exploratory computations of model-call cost per avoided failure or per successful repair by combining provider billing-cost metadata (if available externally) with the archived hosted\_model\_call counts and the archived contingency table; we label any such cost/latency calculations as exploratory because they were not the preregistered primary estimand.

Archival pointers for independent reproduction of reported numbers. Primary analysis artifacts are archived at: analysis/paired\_contingency\_table.csv, analysis/condition\_summary.csv, analysis/missingness\_summary.csv. Raw per-episode records and scoring traces are archived at execution/raw\_results.jsonl and execution/scoring/*.json. The transformation artifact implementing the harmful-change mapping is referenced in execution/transformation\_manifest.json and the full artifact hash manifest is available at execution/execution\_manifest.json (paths recorded in the execution manifest). The execution manifest itself records the controlled\_fault\_assignment\_sha256 = 8d087987\allowbreak{}2348ae94\allowbreak{}a45f3d67\allowbreak{}4df57d65\allowbreak{}52624cd8\allowbreak{}41b62ab2\allowbreak{}39a28034\allowbreak{}7f5275c6 and the preregistration SHA-256 = 0ba66573\allowbreak{}928647f0\allowbreak{}08b2cd5b\allowbreak{}9dde6708\allowbreak{}a61ae86e\allowbreak{}ce55b693\allowbreak{}ace534e0\allowbreak{}1e2b132c; these entries enable machine-verifiable retrieval of the exact artifacts used for the reported analysis.

\section{Discussion}
Interpretation. Under the preregistered controlled-challenge and shared-candidate pairing semantics, exposing deterministic verifier diagnostics substantially increased verifier+simulator-validated configuration success in our synthetic vendor-CLI task bank (approx. +20.51 percentage points). The shared-candidate paired design isolates the causal effect of diagnostic exposure because each arm received the same controlled-fault candidate, identical repair budgets, and identical model identity.

Boundaries and operational implications. (1) Confirmatory scope --- The result applies to the preregistered synthetic controlled-fault bank, the declared model (gpt-5-mini), and the deterministic verifier+simulator used here; external generalization requires multi-model and multi-verifier replications. (2) Safety verdict unresolved --- Because the archived confirmatory summary did not include the preregistered harmful-overrepair aggregation, we make no confirmatory safety verdict on the preregistered <=5\% harmful-overrepair threshold; the archived per-episode traces and transformation manifest permit independent reconstruction. (3) Verifier tax and deployment feasibility --- Hosted-model call counts and stage-level instrumentation are archived; quantifying large-scale operational cost/latency tradeoffs (the verifier tax) requires larger deployments and are outside the confirmatory scope. (4) Verifier-exploitation risk --- Single-verifier iterative loops can be gamed in longer-horizon agentic pipelines; mitigation strategies include multi-verifier ensembles, calibrated abstention, and uncertainty quantification.

Reproducibility and auditability. All artifacts produced by the execution and analysis are archived in the evidence bundle. Key artifact locations: execution/raw\_results.jsonl (raw per-episode records; SHA-256 recorded in execution manifest), execution/scoring/*.json (per-episode scoring traces), analysis/paired\_contingency\_table.csv and analysis/condition\_summary.csv (archived analysis outputs), execution/model\_configuration.json (declared model settings; temperature = null), and execution/transformation\_manifest.json (transformation artifacts including the harmful-change mapping). The execution manifest contains the full hash manifest path (execution/execution\_manifest.json) and representative artifact hashes cited in this manuscript. The exact initial master prompt is archived at provenance/master\_prompt.txt (SHA-256: f781dd79\allowbreak{}78328198\allowbreak{}146d586c\allowbreak{}f33e0f4b\allowbreak{}783c8db1\allowbreak{}26bd0f16\allowbreak{}fcd13699\allowbreak{}455ebb10) and is referenced in the Disclosure Statement.

Threats to validity. Internal validity --- prompt sensitivity, sampling nondeterminism (temperature = null/unset), and the shared\_controlled\_fault\_candidate semantics constrain sampling interpretations: exactly one sampled initial generation per pair was performed and reused; differences between arms arise from diagnostic exposure and subsequent repair calls. Construct validity --- primary success required both deterministic validator and simulator agreement; simulator fidelity limits construct validity. External validity --- single model and synthetic task-bank limit generality to production heterogeneity. Conclusion validity --- n = 39 complete pairs yields detectable effects of the observed magnitude but limits precision for fault-class subgroup inference. These threats are documented and linked to archived artifacts for transparent audit.

\section{Limitations}
\begin{itemize}
\item Synthetic controlled-fault task bank: tasks were deterministically injected and do not represent a broad naturalistic distribution of production NetOps incidents; external validity to live operator workloads is limited.
\item Single-model and single-adapter evaluation: confirmatory execution used only the declared model gpt-5-mini and one adapter family; results do not establish multi-model generality.
\item Simulator-backed primary success predicate: primary success required agreement of deterministic verifier and simulator; simulator fidelity and verifier coverage constrain construct validity.
\item Sample size and power: complete\_pairs = 39 provides detectable effects of the observed magnitude but limits precision for subgroup and per-fault-class inference; exploratory subgroup analyses are not confirmatory.
\item Operational cost/latency unresolved for deployment: execution was instrumented and costs recorded, but large-scale operational feasibility (verifier tax tradeoffs) requires larger-scale study.
\item Verifier-exploitation and long-horizon agentic pathologies remain possible: this controlled setting mitigates some risks, but broader agentic workflows need additional defenses (multi-verifier designs, calibrated abstention, uncertainty quantification).
\item Preregistered confirmatory harmful-overrepair test not present in archived confirmatory summary: the archived confirmatory analysis executor did not compute the preregistered harmful-overrepair aggregation; the per-episode traces and transformation manifest required to compute it are archived and available for independent reaggregation, but the preregistered confirmatory safety verdict is therefore unresolved in this paper.
\end{itemize}

\section{Conclusion}
In a preregistered controlled-challenge paired experiment using hosted\_netops\_validator\_feedback\_repair\_v2 and shared controlled-fault candidates, exposing deterministic verifier diagnostics to an LLM repair workflow produced a statistically detectable and substantial increase in verifier+simulator-validated configuration success (approx. +20.51 percentage points; p = 0.021484375; 95\% CI [0.07692307692307693, 0.358974358974359]). This finding is bounded to the preregistered synthetic task bank, the declared model (gpt-5-mini), and the archived deterministic verifier and simulator. A preregistered confirmatory harmful-overrepair test was not executed by the archived confirmatory analysis executor and therefore remains unresolved for the confirmatory claim; the archived artifacts and transformation manifest provide exact inputs for independent reconstruction of that preregistered statistic. We release the artifact bundle for reproducible reaggregation and recommend multi-model, multi-verifier replications and larger-scale cost/latency studies to assess operational feasibility and safety at NetOps scale.

\section*{Disclosure Statement}
Preregistration: vCLI\_fault\_injection\_repair\_2026\_A\_prereg\_v1 (preregistration\_sha256 recorded in execution manifest). Adapter: hosted\_netops\_validator\_feedback\_repair\_v2. Declared model used for confirmatory execution: gpt-5-mini (declared\_model\_version recorded in execution/model\_configuration.json). Exact initial master prompt archived at provenance/master\_prompt.txt (SHA-256: f781dd79\allowbreak{}78328198\allowbreak{}146d586c\allowbreak{}f33e0f4b\allowbreak{}783c8db1\allowbreak{}26bd0f16\allowbreak{}fcd13699\allowbreak{}455ebb10) and cited here by immutable archive reference. Human role: a Research Director selected the topic family and pre-lock constraints; after the autonomy lock the autonomous system performed literature verification, preregistration, experiment execution, analysis, and manuscript generation; no manual human editing of technical manuscript content occurred after the lock. Execution semantics: execution manifest records temperature = null/unset in execution/model\_configuration.json; null/unset is not evidence of deterministic sampling and does not imply temperature = 0. Pairing semantics: exactly one sampled initial generation per task pair was performed and reused by both arms under shared\_controlled\_fault\_candidate semantics; baseline denotes blind repair (no validator diagnostics exposed) and guarded denotes repair with deterministic validator diagnostics exposed. Harmful-change mapping and reconstructability: the preregistered harmful-change rules were frozen and the transformation artifact implementing the mapping is archived (see execution/transformation\_manifest.json and the full artifact manifest execution/execution\_manifest.json in the evidence bundle). Per-episode raw traces and scoring artifacts required to reproduce all aggregated confirmatory and preregistered secondary statistics are archived (execution/raw\_results.jsonl; execution/scoring/*.json; analysis/*.csv). The study followed the preregistered stopping rule; no silent adapter substitution occurred.

\begin{thebibliography}{99}
\bibitem{ref1} Rafael Cadenas, "Convergent Correctness in Stochastic Code Generation: A Generate-Verify-Repair Architecture for Deterministic Validation of Autonomous AI Coding Agents in Regulated Environments", 2026, 10.2139/ssrn.6754899
\bibitem{ref2} https://arxiv.org/abs/2607.20478
\bibitem{ref3} Arnav Gupta, "Verifier Exploitation in NLI-Guided Iterative Refinement: A Controlled Empirical Analysis", 2026, 10.21203/rs.3.rs-10187492/v1
\bibitem{ref4} http://arxiv.org/abs/2603.19328
\bibitem{ref5} Muhammad Bilal, Jon Crowcroft, Ruizhi Wang, Xiaolong Xu, Schahram Dustdar, "Large Language Models for Agentic NetOps and AIOps: Architectures, Evaluation, and Safety", 2026, 10.48550/arxiv.2605.12729
\end{thebibliography}

\end{document}
